Mục tiêu cốt lõi của danh mục đầu tư rủi ro tối thiểu (MVP) là tìm kiếm một tập hợp các tỷ trọng tài sản $w$ sao cho tổng thể của danh mục được đo bằng phương sai $\sigma_w^2$ là nhỏ nhất, do nhà đầu tư coi rủi ro là điều đáng ghét và lợi nhuận kỳ vọng là điều mong muốn. MVP đại diện cho điểm mà tại đó sự lo ngại rủi ro được đẩy lên mức cao nhất và không quan tâm đến lợi nhuận đạt được bao nhiêu. Để xác định bài toán tối ưu toán học cho danh mục gồm $N$ tài sản rủi ro được xác định bởi vectơ $W = (w_1, w_2, \dots, w_n)$, ta sử dụng ma trận hiệp phương sai $C = [\sigma_{ij}]$ với $\sigma_{ij} = \text{Cov}(R_i, R_j)$ và vectơ suất lợi nhuận $\mu = (\mu_1, \mu_2, \dots, \mu_n)$. Nhà đầu tư giải bài toán tối ưu hóa bằng phương pháp nhân tử Lagrange với hàm mục tiêu là giảm thiểu phương sai danh mục $\min \sigma_w^2 = w^T C w$ (trong đó $w^T$ là ma trận chuyển vị của $w$), đi kèm ràng buộc tổng tỷ trọng các tài sản phải bằng $1$, nghĩa là $\sum_{i=1}^{N} w_i = 1$ nhằm đảm bảo toàn bộ tài sản được đầu tư hết.\\
Ý nghĩa của hiệp phương sai được thể hiện qua công thức:
$$\sigma_w^2 = \sum_{i=1} w_i^2 \sigma_i^2 + 2\sum_{i<j} w_i w_j \sigma_{ij}$$
Công thức này được tính dựa trên trọng số và các phần tử trong ma trận hiệp phương sai. Vai trò của hiệp phương sai giữa các tài sản là rất lớn, nếu hiệp phương sai âm hoặc thấp thì rủi ro tổng thể của danh mục sẽ giảm đáng kể nhờ các biến động ngược chiều triệt tiêu lẫn nhau \cite{macquarie2024exercise}. Khi chuyển sang mô hình Copula, mối liên hệ này sẽ cho thấy rủi ro thực tế không chỉ nằm ở hiệp phương sai mà còn ở cấu trúc phụ thuộc phi tuyến tại các điểm cực trị của đuôi phân phối \cite{mcneil2015quantitative}.\\
Đường biên hiệu quả của mô hình mang lại những kết quả quan trọng về mặt hình học và ứng dụng. Trên mặt phẳng rủi ro - lợi nhuận với trục hoành là độ lệch chuẩn $\sigma$ và trục tung là lợi nhuận kỳ vọng $\mu$, các danh mục khả thi tạo thành một vùng được bao bọc bởi một đường cong hình hyperbol. Điểm MVP của danh mục này nằm chính xác tại đỉnh của đường cong hyperbol, vốn là điểm cực trái nhất trên đồ thị. MVP cũng chính là điểm khởi đầu của đường biên hiệu quả; mọi danh mục nằm phía trên điểm MVP trên đường cong này được coi là hiệu quả, trong khi các danh mục nằm phía dưới MVP là không hiệu quả vì có rủi ro cao hơn nhưng lợi nhuận lại thấp hơn \cite{arratia2014computational}. Sự ra đời của đường biên hiệu quả đã thúc đẩy sự phát triển bùng nổ của các phương pháp quản trị rủi ro trong những thập kỷ tiếp theo, bao gồm các mô hình quan trọng như tỷ lệ Sharpe, mô hình định giá vốn tài sản (CAPM) và lý thuyết kinh doanh chênh lệch giá (APT).